% Formalización pública del resultado autoral de polarización causal.

\section{Champ de polarisation causale et perte de liberté radiale}
\label{sec:md-polarizacion-causal-definicion}

Soit \((\mathcal M,g)\) une géométrie à symétrie sphérique écrite en
coordonnées de Painlevé--Gullstrand,
\begin{equation}
 \mathrm ds^2
 =-c^2\,\mathrm dt^2
  +\bigl(\mathrm dr+v(r)\,\mathrm dt\bigr)^2
  +r^2\,\mathrm d\Omega_2^2,
 \label{eq:md-polarizacion-metrica-pg}
\end{equation}
où \(v(r)\geq0\) mesure l'inclinaison radiale des cônes causaux par rapport
à la congruence temporelle sélectionnée. Le champ adimensionnel
\begin{equation}
 \mathcal P(r):=\frac{v(r)}{c}
 \label{eq:md-polarizacion-parametro}
\end{equation}
est l'intensité de polarisation causale. Sur la distribution spatiale
orthogonale à la congruence temporelle, on définit le tenseur
\begin{equation}
 \Pi(r)=\mathcal P(r)^2\,\widehat r\otimes\widehat r,
 \label{eq:md-polarizacion-tensor}
\end{equation}
dont la direction principale est radiale et dont les directions tangentielles
restent dégénérées. Cette définition n'introduit pas de direction absolue :
la direction principale est déterminée de manière covariante par le gradient du
rayon d'aire.

Les génératrices radiales nulles de \eqref{eq:md-polarizacion-metrica-pg}
satisfont
\begin{equation}
 \frac{\mathrm dr}{\mathrm dt}=-v(r)\pm c
 =c\bigl(-\mathcal P(r)\pm1\bigr).
 \label{eq:md-polarizacion-generatrices}
\end{equation}
Par conséquent, l'ensemble des signes radialement accessibles change de cardinal
lors du franchissement de \(\mathcal P=1\).

\begin{theorem}[Direction causale de chute imposée]
\label{thm:md-polarizacion-caida-obligada}
Dans la géométrie \eqref{eq:md-polarizacion-metrica-pg}, les
trois alternatives exhaustives suivantes sont vérifiées :
\begin{enumerate}
 \item si \(0\leq\mathcal P<1\), une génératrice radiale future augmente \(r\)
 et l'autre le diminue ;
 \item si \(\mathcal P=1\), la génératrice sortante satisfait
 \(\mathrm dr/\mathrm dt=0\) et définit l'horizon ;
 \item si \(\mathcal P>1\), les deux génératrices radiales futures satisfont
 \(\mathrm dr/\mathrm dt<0\).
\end{enumerate}
Dans le troisième régime, le degré radial ne disparaît pas de la variété, mais cesse
de constituer un choix spatial bidirectionnel : tout vecteur causal futur
possède une composante dirigée vers des rayons d'aire plus petits. L'intérieur est, en
ce sens exact, une région de chute causale imposée.
\end{theorem}

\begin{proof}
Les trois conclusions s'obtiennent en comparant les deux valeurs
\(-\mathcal P\pm1\) de \eqref{eq:md-polarizacion-generatrices} à zéro. Pour
\(\mathcal P>1\), même la branche de signe positif est négative. La
conclusion concerne le cône causal et non un choix particulier de
coordonnée temporelle.
\end{proof}

\section{Réalisation de Schwarzschild et gradient collectif de chute}
\label{sec:md-polarizacion-schwarzschild}

Pour une masse totale \(M\), la réalisation de Schwarzschild dans la carte
\eqref{eq:md-polarizacion-metrica-pg} est donnée par
\begin{equation}
 v(r)=c\sqrt{\frac{R}{r}},
 \qquad
 R=\frac{2GM}{c^2},
 \qquad
 \mathcal P(r)^2=\frac{R}{r}.
 \label{eq:md-polarizacion-schwarzschild}
\end{equation}
L'horizon est caractérisé simultanément par
\begin{equation}
 r=R,
 \qquad
 \mathcal P=1,
 \qquad
 g^{\mu\nu}\partial_\mu r\,\partial_\nu r=0.
 \label{eq:md-polarizacion-horizonte}
\end{equation}
L'accélération radiale de la congruence de chute se retrouve comme gradient
de l'intensité cinétique du champ :
\begin{equation}
 a_r=-\frac12\frac{\mathrm d}{\mathrm dr}v(r)^2
 =-\frac{GM}{r^2}.
 \label{eq:md-polarizacion-gradiente}
\end{equation}
Ainsi, la gravitation radiale s'exprime comme gradient collectif de la chute et
non comme une force ajoutée après avoir fixé les trajectoires.

\begin{corollary}[Macropolarisation causale de l'intérieur gravitationnel]
\label{cor:md-macropolarizacion-interior}
Dans la réalisation \eqref{eq:md-polarizacion-schwarzschild}, la région
\(r<R\) satisfait \(\mathcal P>1\). Le trou noir constitue une
macropolarisation causale : la direction principale de \(\Pi\) dépasse le seuil
unitaire et aligne les deux familles radiales futures vers la diminution du
rayon d'aire.
\end{corollary}

\section{Correspondance avec les \texorpdfstring{\(3\)}{3} régimes du TRIT}
\label{sec:md-polarizacion-trit}

Le lecteur géométrique du TRIT se réalise par
\begin{equation}
 \tau(r)=\operatorname{sgn}\!\bigl(1-\mathcal P(r)^2\bigr)
 \in\{+1,0,-1\}.
 \label{eq:md-polarizacion-lector-trit}
\end{equation}
La correspondance est
\begin{equation}
 \begin{array}{c|c|c|c}
 \tau & \mathcal P & \text{régime géométrique} &
 \text{orientation radiale future}\\ \hline
 +1 & \mathcal P<1 & \text{extérieur} & \text{entrée et sortie}\\
 0  & \mathcal P=1 & \text{horizon} & \text{génératrice nulle sortante}\\
 -1 & \mathcal P>1 & \text{intérieur} & \text{descente nécessaire}
 \end{array}
 \label{eq:md-polarizacion-tabla-trit}
\end{equation}
L'involution temporelle du TRIT échange les régimes extérieur et intérieur
et fixe l'horizon. Le changement \(+1\leftrightarrow-1\) n'élimine pas la mémoire
de la route : il modifie son orientation causale et conserve la frontière
\(\tau=0\) comme surface de raccordement.

\Needspace{30\baselineskip}
\section{Clôture locale de matière et alignement gravitationnel collectif}
\label{sec:md-polarizacion-materia}

Soit \(\gamma_a\) une section matérielle persistante de projecteur radial
\(\widehat r_a\otimes\widehat r_a\) et de poids de clôture \(w_a\geq0\). Pour un
ensemble fini de sections, on définit le paramètre d'ordre
\begin{equation}
 \Pi_N
 =\frac{1}{\sum_{a=1}^{N}w_a}
  \sum_{a=1}^{N}w_a\,
  \widehat r_a\otimes\widehat r_a.
 \label{eq:md-polarizacion-colectiva}
\end{equation}
Une particule individuelle réalise une clôture locale persistante de route. L'
intérieur gravitationnel apparaît lorsque le champ collectif fixe une direction
principale causale pour toutes les routes futures. C'est pourquoi on n'identifie pas deux
objets mathématiquement distincts :
\begin{equation}
\begin{aligned}
 \text{particule}
 &=\text{section matérielle persistante},\\
 \text{trou noir}
 &=\text{alignement causal macroscopique au-dessus du seuil}.
\end{aligned}
 \label{eq:md-polarizacion-distincion}
\end{equation}

La séquence causale complète de cette résidence est fixée par
\begin{equation}
\begin{aligned}
 \text{état APP--TRIT--TPK}
 &\longrightarrow
 \text{section matérielle}
 \longrightarrow
 \text{champ de polarisation }\Pi\\
 &\longrightarrow
 \mathcal P=1
 \longrightarrow
 \text{horizon}
 \longrightarrow
 \mathcal P>1
 \longrightarrow
 \text{descente causale nécessaire}.
\end{aligned}
 \label{eq:md-polarizacion-cadena}
\end{equation}
Le chapitre suivant prolonge cette structure vers l'inversion de feuillet,
la purification avancée--retardée, la géométrie de Kruskal et l'information
d'horizon.

\section{Conditions géométriques de la macropolarisation causale : rayon d'aire, inclinaison des cônes, horizon et chute intérieure}
\label{sec:md-polarizacion-control}

Le résultat exige simultanément que le rayon utilisé soit un rayon d'aire, que
\(v(r)\) incline les cônes selon \eqref{eq:md-polarizacion-metrica-pg}, que
l'horizon satisfasse \(\mathcal P=1\) et que les deux pentes de
\eqref{eq:md-polarizacion-generatrices} soient négatives à l'intérieur. Ces
identités distinguent la perte effective de liberté radiale d'une simple
métaphore dimensionnelle et fournissent un contrôle direct de la
macropolarisation causale.
